% We propose a new frontier: neural computers (NCs)—learned, executable runtimes—toward “living computers” that learn interface dynamics from experience.  
% We call the long-term goal a Completely Neural Computer (CNC): a single learned model of a computer that subsumes the classical compute/memory/I/O stack, thereby challenging the layered OS/application architecture that mediates I/O.
% As an initial step toward this vision, we ask whether a neural model can learn computer interface dynamics from collected I/O samples.
% We instantiate NCs as video models that roll out screen frames from instructions, recent pixels, and (when available) user actions, without access to instrumented program state.
% We conduct experiments in both CLI and GUI settings and derive practical design considerations for building CNCs.
% Finally, rather than imitating human or animal cognition, we ask whether neural networks can be designed to mirror the computational structure of machines—like circuit-based computers—and thereby go beyond both current foundation models and conventional computers.



We propose a new frontier: Neural Computers (NCs)—neural networks that unify computation, memory, and I/O in a single latent runtime state.
Our long-term goal is the Completely~Neural~Computer (CNC), the fully realized form of this vision: a general-purpose neural computer that challenges the layered architecture of conventional computers.
As an initial step, we study whether such executable dynamics can be learned~solely~from collected I/O traces, without access to instrumented program state.
Concretely, we instantiate NCs as video models that roll out screen frames conditioned on instructions, pixels, and (when available) user actions, and evaluate them in both CLI~and~GUI~settings.
We outline a roadmap toward CNCs, centered on fundamental challenges. 
If achieved, CNCs could inaugurate a new computing paradigm beyond both today’s foundation models and conventional~computers.
